\scanpage{27}
% archive foliation: 26
\diagram{\begin{tikzcd}[column sep=large]
\operatorname{Hom}(\Psi^i_{\deleted{\unsure{*}}},\Phi(j,n))\arrow[r]&\operatorname{Hom}(\Psi^i,\Psi^j)\arrow[d,phantom,"=" description]\\
&0\text{ si }j\neq i
\end{tikzcd}}
\unsure{$\forall\alpha\in$}
donc si $j\neq i$, $\operatorname{Hom}(\Psi^i,\Phi(j,n))\xleftarrow{\sim}$
$\operatorname{Hom}(\Psi^i,\Psi^{j-1}\deleted{\unsure{\mathrm{can}}})$,
\deleted{donc}
qui est nul si
$i\neq j-1$ i.e. $j\neq i+1$; si $j=i+1$, on trouve $\simeq k$.
Pour $j=i$, la composée
$\Psi^i\xrightarrow{\alpha}\Phi(i,n)\xrightarrow{\mathrm{can}}\Psi^i$
est de la forme
$d\operatorname{id}_{\Psi^i}$, donc l'image
inverse de l'extension
\diagram{\begin{tikzcd}[column sep=large]
&&\Psi^i\arrow[d,"\alpha"']\arrow[dr]\\
0\arrow[r]&\Psi^{j-1}\arrow[r,dash]&\Phi(j,n)\arrow[r,dash]&\Psi^i\arrow[r,dash]&0
\end{tikzcd}}
$\Phi(i,n)$ est de classe $d e_{i,n}=dn e_i$,
qui est nulle, donc $dn e_i=0$, donc
$dn=0$, donc $d=0$, donc $\unsure{\alpha}$ se factorise
par $\Psi^i\longrightarrow\Psi^{i-1}$, donc est nul.
\deleted{Pour $j=i+1$, on a $\operatorname{Hom}($}
Donc
\deleted{$\unsure{\operatorname{Hom}(\Phi,\Psi)}$}
\[
\operatorname{Hom}(\Psi^i,\Phi(j,n))\simeq
\left\{\begin{array}{ll}
0&\text{si }j\neq i+1\\
k&\text{si }j=i+1\quad\text{(par l'hom can. }\Psi^i\longrightarrow\Phi(i+1,n)\text{)}
\end{array}\right.
\]
\par\bigskip\hrule\bigskip
Enfin on suppose
\[
\text{(d)}\qquad\begin{gathered}
\text{Tout objet de }C\text{ est somme}\\
\text{finie d'objets de la forme}\\
\Phi(i,n)\ \inserted{(i\geq1)}\text{ ou }\Psi^j\ (j\geq0)
\end{gathered}
\]
